% =============================================================================
%  Chalmers Beamer template -- self-documenting tutorial deck
%
%  Unofficial Beamer theme approximating the 2026 Chalmers visual identity.
%  Vibe-coded by Wenliang Zhang; not an official Chalmers product.
%
%  Project / issue tracker:
%      https://github.com/wenlzhang/chalmers-beamer
%
%  Licence: MIT for the LaTeX/Beamer code; the Chalmers logotypes in
%           logographics/ remain Chalmers trademarks (see LICENSE +
%           NOTICE in the repository root).
%
%  This file doubles as (a) a minimal compilable example you can copy as a
%  starting point and (b) a slide-by-slide tutorial of every feature the
%  Chalmers Beamer theme provides.  Read the rendered PDF top-to-bottom for
%  the full guide; read this .tex file for the underlying LaTeX patterns.
%
%  Compile:   pdflatex your-deck.tex   (run twice -- once for the .toc,
%                                       once to render the section dividers)
%
%  See README.md for an out-of-deck reference of every command.
% =============================================================================

% --- Aspect ratio (in \documentclass options) -------------------------------
% Common Beamer aspect ratios (set via aspectratio=NNN):
%   169  -- 16:9   (modern widescreen, matches the Chalmers PPT master)
%   1610 -- 16:10  (older widescreen, MacBook displays)
%   149  -- 14:9
%   54   -- 5:4
%   43   -- 4:3    (Beamer's default; older projectors)
%   32   -- 3:2
% Recommended: aspectratio=169 (the official Chalmers PPT template is 16:9).
%
% --- Handout / notes mode (optional) ----------------------------------------
%   handout   -- one slide per page, all overlays collapsed
%   trans     -- transparency mode (no overlays)
%   notes=...  -- speaker-note rendering (see Beamer manual)
\documentclass[aspectratio=169]{beamer}
% \documentclass[aspectratio=169,handout]{beamer}
% \documentclass[aspectratio=43]{beamer}

\usetheme{Chalmers}

% --- Common packages --------------------------------------------------------
\usepackage[utf8]{inputenc}
\usepackage[T1]{fontenc}
\usepackage[british]{babel}
\usepackage{amsmath,amssymb}
\usepackage{siunitx}
\usepackage{graphicx}
\usepackage{booktabs}
\usepackage{caption,subcaption}
\usepackage{tikz}
\usepackage{pgfplots}
\pgfplotsset{compat=newest}

\usepackage{hyperref}
\hypersetup{colorlinks=true,allcolors=ChalmersPurple}

% --- Speaker notes -----------------------------------------------------------
% \setbeameroption{show notes}                          % notes inline
% \setbeameroption{show notes on second screen=right}   % presenter view
\setbeameroption{hide notes}

% =============================================================================
%  Metadata  (used by \titlepage and the footer)
% =============================================================================
\title{Chalmers Beamer Template}
\subtitle{A self-documenting tutorial deck}
\shorttitle{Chalmers Beamer Tutorial}

\author{Wenliang Zhang}
%\role{Role}
\institute{Chalmers University of Technology}
\date{2026-05-22}
%\place{@Place}
% --- Logo style on every slide (optional) -----------------------------------
% Default `logo` (avancez emblem only).  Uncomment ONE to switch:
%   \chalmersheadline{logo}           % avancez emblem only          (DEFAULT)
%   \chalmersheadline{word}           % CHALMERS wordmark only
%   \chalmersheadline{logo+word}      % avancez + CHALMERS (compact)
%   \chalmersheadline{logo+word+sub}  % avancez + CHALMERS + UNIV. OF TECH.

% --- Section-divider style (optional) ---------------------------------------
%   \chalmerssectiondividers{progress}   % all sections, current highlighted (DEFAULT)
%   \chalmerssectiondividers{solid}      % purple full-bleed with section name
%   \chalmerssectiondividers{none}       % no dividers (short decks)

% --- Outline frame (optional) ------------------------------------------------
%   \chalmersoutline{show}               % \outlineframe renders the TOC  (DEFAULT)
%   \chalmersoutline{hide}               % \outlineframe is a no-op (drop the outline)

% --- Page-number format in the footer (optional) -----------------------------
%   \chalmersfooterpages{current}        % "5"        (DEFAULT)
%   \chalmersfooterpages{both}           % "5 / 20"   (current / total)

% =============================================================================
\begin{document}

% --- Cover slide ------------------------------------------------------------
\titleframe

% --- Outline ----------------------------------------------------------------
% Hide the outline entirely by adding \chalmersoutline{hide} in the preamble.
\outlineframe

% =============================================================================
\section{Quick start}
% Covers: what the template provides, minimal source, file layout,
%         the three slide types, and frame alignment options.
% =============================================================================

\begin{frame}{What this template provides}
    \footnotesize
    \begin{alertblock}{Unofficial template}
        This is an \emph{unofficial} community Beamer theme that
        approximates the 2026 Chalmers visual identity using the Chalmers logos and the official PowerPoint master as reference.
        It is not an official Chalmers product.  For brand-critical
        materials consult the Chalmers Brand Portal.
    \end{alertblock}

    \medskip
    \normalsize

    \begin{itemize}
        \item Primary brand colour: \textbf{Chalmers Purple} (\texttt{\#472DBE})
        \item Full accent palette (green, cyan, red, orange, yellow, pink, \ldots)
        \item Three slide types: cover, section divider, normal frame
        \item Four logo styles, three section-divider styles -- pick once,
              applies everywhere
        \item Arial-like typography (Helvetica via the \texttt{helvet} package)
    \end{itemize}
\end{frame}

\begin{frame}[fragile=singleslide]{Minimal source}
    \scriptsize
    \begin{verbatim}
\documentclass[aspectratio=169]{beamer}
\usetheme{Chalmers}

\title{...}
\subtitle{...}
\shorttitle{Short title for footer}
\author{...}
\institute{...}
\date{2026-05-22}

\begin{document}
    \titleframe                        % cover slide
    \begin{frame}{Outline}\tableofcontents\end{frame}
    \section{Introduction}             % auto-inserts a divider
    \begin{frame}{Topic}Hello, world.\end{frame}
\end{document}
    \end{verbatim}
\end{frame}

\begin{frame}{File layout}
    \footnotesize
    \begin{tabular}{@{}lp{0.55\textwidth}@{}}
        \toprule
        File & Purpose \\
        \midrule
        \texttt{beamerthemeChalmers.sty} & The theme (colours, fonts, layout) \\
        \texttt{logographics/}           & Chalmers logos (PNG) \\
        \texttt{your-deck.tex}           & This skeleton / tutorial deck \\
        \texttt{README.md}               & Full command reference \\
        \bottomrule
    \end{tabular}

    \medskip

    Drop the four items into any folder, edit \texttt{your-deck.tex},
    run \texttt{pdflatex your-deck.tex} twice (the second pass renders
    the section dividers from the \texttt{.toc} file).
\end{frame}

% =============================================================================
% (Slide-types content folded into Quick start)
% =============================================================================

\begin{frame}{Three built-in slide types}
    \begin{itemize}
        \item \textbf{Cover slide} -- purple background, white logo
              top-right, title / subtitle / author / date.\\
              Invoke with \texttt{\textbackslash titleframe}.
        \item \textbf{Section divider} -- purple background, list of all
              sections with the current one highlighted.\\
              Inserted automatically by every \texttt{\textbackslash
              section\{\ldots\}}.
        \item \textbf{Normal content frame} -- white background, purple
              title, logo top-right, footer with date $|$ short title $|$
              slide number.
    \end{itemize}

    \medskip

    For custom layouts on a purple background (e.g.\ a closing splash,
    a hand-crafted divider) use the \texttt{chalmerscover} environment.
\end{frame}

\begin{frame}[fragile]{Custom purple frame -- \texttt{chalmerscover} env}
    \small
    \begin{verbatim}
\begin{chalmerscover}
    \vfill
    \begin{center}
        {\color{white}\Huge\bfseries Thank you}\\[6mm]
        {\color{white}\Large Questions?}
    \end{center}
    \vfill
\end{chalmerscover}
    \end{verbatim}

    Inside the environment you automatically get: purple background, white
    logo in the top-right corner (same position as every other slide), no
    footer.  Drop in any content.
\end{frame}

\begin{frame}[t]{Frame alignment options}
    Beamer's standard \texttt{[t]}, \texttt{[c]}, \texttt{[b]} options
    work normally with this theme:

    \medskip

    \begin{tabular}{@{}lp{0.65\textwidth}@{}}
        \toprule
        Option & Effect \\
        \midrule
        \texttt{[t]} & content anchored to the top of the body area \\
        \texttt{[c]} & content vertically centred (Beamer default) \\
        \texttt{[b]} & content anchored to the bottom \\
        \bottomrule
    \end{tabular}

    \medskip

    \textit{This very slide uses} \texttt{[t]} \textit{-- the text sits
    just below the title.}
\end{frame}

% =============================================================================
\section{Visual identity}
% Covers: logos, colour palette, layout building blocks (blocks,
%         equations, tables), and section-divider presets.
% =============================================================================

\begin{frame}[fragile]{Logo style switch}
    \texttt{\textbackslash chalmersheadline\{\textit{key}\}} selects which
    mark appears in the top-right corner of every slide:

    \medskip

    \begin{tabular}{@{}lp{0.65\textwidth}@{}}
        \toprule
        Key & Logo \\
        \midrule
        \texttt{logo}            & avancez emblem only \quad\textbf{(default)} \\
        \texttt{word}            & CHALMERS wordmark only \\
        \texttt{logo+word}       & avancez emblem + CHALMERS \\
        \texttt{logo+word+sub}   & emblem + CHALMERS + UNIVERSITY OF TECHNOLOGY \\
        \bottomrule
    \end{tabular}

    \medskip

    The colour adapts automatically: purple on light slides, white on the
    cover and section dividers.
\end{frame}

\begin{frame}{Colour palette -- primary and neutrals}
    \footnotesize
    \begin{tabular}{@{}llp{0.45\textwidth}@{}}
        \toprule
        Name & Hex & Role \\
        \midrule
        \texttt{ChalmersPurple}      & \#472DBE & primary; titles, bullets, rules \\
        \texttt{ChalmersPurpleDark}  & \#332085 & dark variant (pressed states) \\
        \texttt{ChalmersPurpleLight} & \#6746EB & lighter purple \\
        \texttt{ChalmersPurpleSoft}  & \#9E92E8 & soft purple (backgrounds) \\
        \texttt{ChalmersCream}       & \#F0EDE6 & warm off-white (block bodies) \\
        \texttt{ChalmersBlack}       & \#222222 & body text \\
        \texttt{ChalmersGrey}        & \#7F7F7F & subtitles, muted text \\
        \bottomrule
    \end{tabular}
\end{frame}

\begin{frame}{Colour palette -- accents (data viz)}
    \footnotesize
    \begin{tabular}{@{}llp{0.45\textwidth}@{}}
        \toprule
        Name & Hex & Suggested use \\
        \midrule
        \texttt{ChalmersGreen}  & \#00B356 & success, example block \\
        \texttt{ChalmersLime}   & \#ADD739 & accent \\
        \texttt{ChalmersCyan}   & \#36B7F6 & accent \\
        \texttt{ChalmersTeal}   & \#61E9D2 & accent \\
        \texttt{ChalmersRed}    & \#D52515 & alert block, warnings \\
        \texttt{ChalmersOrange} & \#ED7122 & accent \\
        \texttt{ChalmersYellow} & \#F6D62A & highlight \\
        \texttt{ChalmersPink}   & \#D987BA & accent \\
        \bottomrule
    \end{tabular}

    \medskip

    All names are defined as \texttt{xcolor} colours -- use them anywhere:
    \texttt{\textbackslash textcolor\{ChalmersGreen\}\{\ldots\}},
    \texttt{fill=ChalmersTeal}, etc.
\end{frame}

% =============================================================================
% (Layout building blocks -- folded into Visual identity)
% =============================================================================

\begin{frame}{Blocks}
    \footnotesize
    \begin{columns}[T,onlytextwidth]
        \begin{column}{0.49\textwidth}
            \begin{block}{Standard block}
                Purple title, cream body. Use for key takeaways.
            \end{block}
            \begin{alertblock}{Alert block}
                Red title, faint red body. Use for warnings.
            \end{alertblock}
            \begin{exampleblock}{Example block}
                Green title, faint green body. Use for examples.
            \end{exampleblock}
        \end{column}
        \begin{column}{0.49\textwidth}
            \textbf{When to use which}
            \begin{itemize}
                \item Plain \texttt{block} -- summary, definition, "key idea"
                \item \texttt{alertblock} -- caveat, exception, common mistake
                \item \texttt{exampleblock} -- concrete instance, demo
            \end{itemize}
        \end{column}
    \end{columns}
\end{frame}

\begin{frame}{Equations and units}
    Numbered equations work as usual:
    \begin{equation}
        \dot{\mathbf{x}} = f(\mathbf{x},\mathbf{u}),
        \label{eq:state}
    \end{equation}
    where $\mathbf{x}\in\mathbb{R}^n$ is the state vector and
    $\mathbf{u}\in\mathbb{R}^m$ is the control input.

    \medskip

    Units use the \texttt{siunitx} package:
    $v_{\max} = 120 \, \operatorname{km/h}$,
    $g = 9.81 \, \operatorname{m/s^2}$.

    \medskip

    Itemize bullets:
    \begin{itemize}
        \item Use Chalmers Purple as the bullet colour automatically.
        \item Sub-bullets use the lighter purple variant.
    \end{itemize}
\end{frame}

\begin{frame}{Tables (booktabs style)}
    \centering
    \begin{tabular}{@{}lcc@{}}
        \toprule
        Quantity & Symbol & Value \\
        \midrule
        Vehicle mass    & $m$        & \SI{1500}{\kilo\gram}        \\
        Wheelbase       & $L$        & \SI{3.0}{\meter}             \\
        Max speed       & $v_{\max}$ & $35 \, \operatorname{m/s}$   \\
        \bottomrule
    \end{tabular}

    \bigskip

    Plain \texttt{tabular} + \texttt{booktabs} (no vertical lines, three
    horizontal rules).  Caption goes \emph{above} the table.
\end{frame}

% =============================================================================
% (Section dividers -- folded into Visual identity)
% =============================================================================

\begin{frame}[fragile]{Three section-divider presets}
    Pick the look of slides inserted between sections via:

    \begin{verbatim}
\chalmerssectiondividers{progress}   % DEFAULT
\chalmerssectiondividers{solid}
\chalmerssectiondividers{none}
    \end{verbatim}

    \begin{tabular}{@{}lp{0.65\textwidth}@{}}
        \toprule
        Preset & What you see \\
        \midrule
        \texttt{progress} & list of all sections, current big + bold + white,
                            others smaller and dim ("you are here") \\
        \texttt{solid}    & purple full-bleed with just the section name \\
        \texttt{none}     & no divider slide at all \\
        \bottomrule
    \end{tabular}
\end{frame}

\begin{frame}{Fully custom dividers}
    For a one-off layout, redefine \texttt{\textbackslash AtBeginSection}
    in your own \texttt{.tex} after \texttt{\textbackslash
    usetheme\{Chalmers\}}.  Inside, build the divider however you like;
    typically wrap content in \texttt{\textbackslash
    begin\{chalmerscover\}\ldots\textbackslash end\{chalmerscover\}} to
    keep the purple-bg + white-logo style.
\end{frame}

% =============================================================================
\section{Customisation}
% =============================================================================

\begin{frame}[fragile]{Override colours and fonts}
    \small
    Place any of these \emph{after} \texttt{\textbackslash
    usetheme\{Chalmers\}}:

    \small
    \begin{verbatim}
% recolour the primary purple
\definecolor{ChalmersPurple}{HTML}{332085}

% recolour the frame title
\setbeamercolor{frametitle}{fg=ChalmersBlack}

% reduce the frame-title size
\setbeamerfont{frametitle}{family=\sffamily,series=\bfseries,size=\large}
    \end{verbatim}

    The theme defines every visible element via \texttt{\textbackslash
    setbeamercolor} / \texttt{\textbackslash setbeamerfont} so any of them
    can be overridden in this way.
\end{frame}

\begin{frame}[fragile]{Override the footer}
    \small
    The default footer is \emph{date $|$ short title $|$ slide number}.
    Override locally by redefining \texttt{footline}:

    \small
    \begin{verbatim}
\setbeamertemplate{footline}{%
    \usebeamercolor[fg]{footline}\usebeamerfont{footline}%
    \vspace{1mm}%
    \hbox to \paperwidth{\hfill\insertframenumber\hspace{8mm}}%
    \vspace{1mm}%
}
    \end{verbatim}

    Available macros inside the template: \texttt{\textbackslash
    insertdate}, \texttt{\textbackslash insertshorttitle@chalmers},
    \texttt{\textbackslash insertframenumber}, \texttt{\textbackslash
    insertauthor}, \texttt{\textbackslash insertinstitute}.
\end{frame}

\begin{frame}[fragile]{Adding a partner logo}
    For collaboration slides, drop the partner's PNG into
    \texttt{logographics/} and place it manually with TikZ:

    \scriptsize
    \begin{verbatim}
\pgfdeclareimage[height=8mm]{partner}{logographics/your-partner-logo}
% inside any frame body:
\begin{tikzpicture}[remember picture,overlay]
    \node[anchor=north east,xshift=-50mm,yshift=-5mm]
        at (current page.north east) {\pgfuseimage{partner}};
\end{tikzpicture}
    \end{verbatim}
    \normalsize

    Tune \texttt{xshift} so the partner logo doesn't collide with the
    Chalmers headline logo in the top-right.
\end{frame}

% =============================================================================
% (Wrap-up -- folded into Customisation as the last two slides)
% =============================================================================

\begin{frame}{Summary}
    \begin{itemize}
        \item \textbf{Identity:} the 2026 Chalmers visual identity, primary
              colour Chalmers Purple (\texttt{\#472DBE}).
        \item \textbf{Three slide types:} cover, section divider, normal
              -- the logo sits at the \emph{same} top-right position on all
              of them.
        \item \textbf{Pick once, applies everywhere:} a single
              \texttt{\textbackslash chalmersheadline} +
              \texttt{\textbackslash chalmerssectiondividers} setting in the
              preamble configures every slide.
        \item \textbf{Customisable:} every colour and font is a Beamer
              setting; override locally without touching the \texttt{.sty}.
    \end{itemize}
\end{frame}

\begin{frame}{Further reading}
    \begin{itemize}
        \item \texttt{README.md} -- full command reference, colour
              palette, customisation examples
        \item \texttt{beamerthemeChalmers.sty} -- the theme itself,
              well commented section-by-section
        \item Chalmers Brand Portal -- official brand guidelines and
              ceremonial-use rules for the avancez emblem
    \end{itemize}
\end{frame}

% --- Closing "Thank you" slide ----------------------------------------------
% Uses the `chalmerscover` environment: same purple-background style as
% section pages, white logo top-right.  Replace the centred block below with
% any custom content.  For a white-background variant, use a normal frame.
\begin{chalmerscover}
    \vfill
    \begin{center}
        {\color{white}\Huge\bfseries Thank you}\\[6mm]
        {\color{white}\Large Questions?}\\[10mm]
        {\color{white!85}\small Bug reports and feature ideas welcome:}\\[2mm]
        {\small\href{https://github.com/wenlzhang/chalmers-beamer}%
            {\color{white}\underline{github.com/wenlzhang/chalmers-beamer}}}
    \end{center}
    \vfill
\end{chalmerscover}

\end{document}
